\chapter*{Resume}
\addcontentsline{toc}{chapter}{Resume}

Ce projet de fin d'études traite la réduction de la surface d'attaque dans un contexte d'administration de systèmes critiques, à travers une architecture Zero Trust intégrée. Le problème principal identifié est l'existence d'accès privilégiés permanents, peu tracés et insuffisamment corrélés au niveau de criticité des actifs.

La solution proposée combine cinq briques principales : \textbf{Authentik} pour la gestion des identités (IAM), \textbf{Tailscale ACL} pour le contrôle d'accès réseau, une \textbf{application web Flask avec PostgreSQL} pour la gouvernance des demandes, un \textbf{moteur Python Just-in-Time} pour l'approbation/révocation automatisées, et \textbf{Wazuh} pour l'observabilité et la traçabilité de sécurité.

Le flux cible impose qu'aucun accès privilégié ne soit permanent : chaque élévation est demandée, justifiée, approuvée, limitée dans le temps (TTL) puis retirée automatiquement. Les scénarios de validation couvrent l'authentification forte, la demande JIT, l'activation ACL, l'accès SSH pendant la fenêtre autorisée, puis le refus après expiration.

Les résultats attendus montrent une amélioration significative de la posture de sécurité : suppression des accès permanents non justifiés, traçabilité complète du cycle d'accès, et meilleure capacité d'audit et de conformité.
En complément, le chapitre sur la détection d'anomalies présente une solution ciblant les comptes administrateurs. Un pipeline non supervisé, basé sur l'algorithme \textit{Isolation Forest}, analyse des indicateurs temporels et comportementaux (heure, jour, indicateur week-end/nuit, nouveau sous-réseau, nombre d'approbations récentes, type d'action, statut) et entraîne un modèle par administrateur ; tout score élevé (seuil opérationnel typique : 0.60) génère une alerte enrichie envoyée au SIEM pour investigation. Les visualisations facilitent l'interprétation des résultats et la traçabilité des détections.

\textbf{Mots-cles :} Zero Trust, IAM, PIM, PAM, Authentik, Tailscale, JIT, ACL, Wazuh, cybersecurite.

\clearpage
